\chapter{跨载体实现关系：碳硅比较、设计闭环与证据边界}

\begin{introduction}
    \item 碳基生命与硅基系统可以比较功能接口，但不能由表面对应推出整体数学同构
    \item 跨载体共同原则来自开放系统、资源预算、结构约束和验证闭环，而非统一物理实体
    \item 生物设计必须经过多尺度模型、实验校准和分级治理，计算生成不能替代现实证据
\end{introduction}

前一章把机器族群的经验更新限定为带来源、版本和回滚的工程事务，本章转向更容易产生越界推断的问题：碳基生命与硅基智能究竟能在什么意义上比较。我们承认两者都可能形成感知、状态保持、调节、学习和行动闭环，也都受资源、时延、噪声和环境扰动限制；但共同使用“信息”“网络”或“目标”等词，不足以证明DNA、神经调节、预训练权重、提示词和优化器属于同一个数学对象。我们将比较单位限定为给定任务、观测协议和干预集合下的实现关系，分别记录材料机制、状态空间、时间尺度、单位、可观测量与失败方式。HSF-HD在此提供跨系统的状态动力学语言，AgentOS只是一组硅基实例；生物实验事实、计算模型结果和工程设想必须分栏保存。

\section{任务相关实现关系：从同构宣称到可检验比较}
\label{sec:isomorphism_map}

设碳基系统在时间窗$[0,T]$内的状态轨迹为$x_C\in\mathcal{X}_C$，其成分可以包含基因调控、代谢物浓度、膜电位和组织环境；硅基系统轨迹为$x_S\in\mathcal{X}_S$，其成分可以包含参数、激活、外部记忆、控制模式和计算资源。两者的坐标、单位与演化算子不同，通常不存在保持所有结构的双射。为了讨论某项功能，我们先给出有限查询族$\mathcal{Q}$、允许干预集$\mathcal{I}$、输出度量$d_q$和容差$\varepsilon_q$，再定义编码$\phi$与读出$R_C^q,R_S^q$。若对声明的初始条件和干预有
\begin{equation}
 d_q\!\left(R_C^q(x_C^{I}),R_S^q(\phi(x_C)^{I})\right)\leq\varepsilon_q,
 \qquad q\in\mathcal{Q},\ I\in\mathcal{I},
\end{equation}
我们只称硅基模型在该协议下近似实现了碳基系统的一组读出。这个定义不要求材料相同，也不承诺未列入查询的性质被保存；扩大时间窗、改变环境或加入新干预后，关系必须重新检验。严格同构还要求映射可逆并保持指定运算，而生物和计算系统常有不可观测状态、有损编码和不可逆事件，因此不应把近似实现升级为“同一物理定律的投影”。

原比较表把基因组对应预训练权重、调控区对应奖励模型或提示词、受体对应Tokenizer、信号通路对应激活、细胞核对应优化器。这些配对可用作提出问题的索引，却不是一一映射。基因组并不独自决定生物世界图，其表达依赖细胞类型、发育史和环境；预训练权重也不是只读基因，部署系统还受检索、工具、提示和运行时状态影响。启动子、增强子、奖励模型与提示词都能调制行为，但前两者参与分子结合和转录动力学，后两者作用于训练目标或条件推断，时间尺度和因果接口不同。细胞膜与传感器都处于边界，却不能把受体信号、模态编码和符号切分视为同一个换能过程。表格因此应增列“可比较功能”“材料机制”“验证干预”和“不保留性质”，让相似性承担启发任务而非证明任务。

局部结构网络与全局分布表示在两类系统中也必须区分。生物调控网络的节点和边来自实验定义，分布活动则可能跨细胞、组织或神经群体；硅基系统的对象图、工具依赖和任务谱系属于结构记录，向量激活与参数表示属于分布状态。我们通过任务身份、来源绑定、时间对齐和读出接口联合它们，不从某个脑区、细胞器或软件模块的名称推出机制等同。例如同一损伤修复结果可能由不同分子通路实现，同一软件行为也可能由不同内部表示产生；结果相似只支持输出层关系，不能反推出内部结构唯一。反过来，内部结构相似也不保证环境结果相同，因为传感器条件、执行器和外部反馈可能改变闭环。

跨载体比较要用对称实验，避免只挑选支持对应关系的例子。我们至少记录初始状态分布、输入扰动、测量分辨率、缺失数据、干预强度、恢复窗口和失败终点；然后分别估计准确性、校准误差、时延、资源成本与风险。秒、摩尔浓度、焦耳、字节、浮点运算和无量纲概率不可直接相加，确需形成综合排序时先声明换算和权重，否则使用词典序或Pareto前沿。预测有效是模型在某个协议下的经验成绩，不证明系统整体同构，更不证明生物具有软件模块或软件获得生命属性。跨载体“映射”最有价值的产物，是一份可反驳的接口假说清单：哪些读出近似保持、哪些干预暴露差异、在哪些尺度关系失效。

\section{共同工程原则：预算、耦合、探索与恢复}
\label{sec:dual_track_engineering}

碳基与硅基系统确有若干共同工程约束，但它们来自开放系统和有限控制的一般事实。第一项是收支闭合。我们为任一载体定义资源向量$b(t)=(b_1,\ldots,b_m)$，每一分量保留自己的单位，输入率$r_{\mathrm{in}}$、维持消耗$r_{\mathrm{keep}}$、任务消耗$r_{\mathrm{task}}$与排出或释放$r_{\mathrm{out}}$满足
\begin{equation}
 \dot b(t)=r_{\mathrm{in}}(t)-r_{\mathrm{keep}}(t)-r_{\mathrm{task}}(t)-r_{\mathrm{out}}(t).
\end{equation}
这只是记账模型，具体守恒关系取决于对象。生物需要物质、化学能和稳态调节，计算系统需要电能、存储、带宽、散热和维护；信息熵变化不能替代焦耳能量平衡。细胞无限分裂、推理无限递归和训练损失趋近零是三种不同问题，不能共同归因为“排熵失败”。工程上真正共享的原则，是为资源设置硬上限、测量饱和、识别外部供给变化并准备降级与恢复。

第二项是结构能力与调节目标的接口化。对载体$p\in\{C,S\}$，令结构状态$s_p$描述较慢的连接、参数或物理构型，调节状态$v_p$描述较快的目标、门控或运行条件，行动由$u_p=\pi_p(x_p;s_p,v_p)$产生。把$s_p$和$v_p$分开能帮助定位错误，却不意味着二者可以独立设计：目标可能超出结构能力，结构更新也会改变目标的可实现集合。羊的摄食与消化例子只能说明能力—需求失配的可能性，不能把欲望等同规范场；RLHF、提示和预训练也在不同阶段作用，价值信息并非只存在于一个附加层。我们要求每次更新给出兼容性测试$K(s_p,v_p)$、硬可行集和回退版本，在目标、结构或环境变化时重算，而不是假定底座永远稳固、调节层天然安全。

第三项是探索压力的可控生成。学习需要尚未解决的问题、数据差异或外部目标，但不需要在系统中植入不可消除的痛苦、死亡恐惧或“拓扑孔洞”。我们把探索驱动拆成不确定度、任务价值、新颖度和风险四个可测量输入，并通过有界调度器选择实验；当潜在伤害超过阈值时，即使新颖度很高也不得执行。目标过易可能减少更新信号，却不必导致系统停止全部计算；目标永久不可达则可能造成资源浪费、欺骗性优化或反复失败。碳基研究需要伦理审查和伤害终点，硅基系统需要沙箱、能力边界和预算终止。两者共享的是受约束实验逻辑，不是“挖洞产生角动量”的物理机制。

最后是恢复优先于完美。生物干预和软件更新都可能出现域外效应，所以设计应保留对照、来源记录、阶段性提交和停止条件。生物系统往往难以真正回滚，必须更重视可逆分子开关、剂量阶梯、隔离与长期随访；硅基系统可以恢复软件版本，但外部动作和已传播数据仍可能不可逆。我们分别测量内部状态恢复、外部后果修复和证据完整性，不把一次目标下降当作成功。HSF-HD可以用快状态、慢结构和离散事件描述这些闭环，但不能由统一记号消除材料差异。真正通用的法则是：先声明对象和单位，再提出有限实现关系；先筛硬约束，再优化软目标；先在小范围验证，再扩大干预；失败时保全证据并进入已演练的恢复路径。

这四项原则还要求负结果能够进入模型。若某个跨载体映射只在成功案例上成立，我们应缩小适用域或撤销对应关系，而不是把失败解释为系统“尚未达到足够高阶”。同样，生物实验的统计显著、模拟器中的适应度提升和软件基准分数上升属于不同证据等级；只有在预先声明的复现实验、外部验证和风险审查中，它们才可能共同支持一个设计决定。统一语言的作用是暴露差异并组织验证，不是让差异消失。

% Draft continues at sec:reverse_evolutionary_computation
